\section{Future work}
\label{sec:futurework}

The three algorithms implemented here exploit three structural properties of
the fBM covariance matrix: positive-definiteness (Cholesky), stationarity of
increments (circulant FFT), and smooth off-diagonal low-rank structure (global
rSVD).

\subsection{True hierarchical H-matrix}
\label{sec:hmatrix-future}

The global rSVD (Algorithm~3) approximates the entire $N \times N$
matrix $C$ at rank $k$.  This forces the rank budget to cover both the
smooth far-off-diagonal region and the singular near-diagonal region, where
$C(t,s) \sim |t-s|^{2H}$ with $H = 0.10$ is Hölder-$0.2$ continuous and
barely above $C^0$ regularity.  The global approximation cannot specialize, which is
why the spectrum of the full matrix decays slowly at small $H$ even though its
off-diagonal block compresses to rank $\leq 3$ (Figure~\ref{fig:structure}(d)).

A true hierarchical H-matrix separates these concerns by a recursive
block-tree decomposition of $C$ \citep{hackbusch2015hierarchical,
bebendorf2008hierarchical}.  At each level of the tree:
\begin{itemize}
  \item Near-diagonal blocks (time intervals that overlap or touch)
        are stored dense.  The near-diagonal singularity is not low-rank, so
        no rank budget is wasted trying to compress it.
  \item Far-off-diagonal blocks (well-separated time intervals) lie
        entirely in the smooth region $t \neq s$ and are approximated at low
        rank proportional to the distance from the diagonal.
\end{itemize}
The result is an H-Cholesky factorization with $O(Nk\log N)$ construction
cost and $O(Nk\log N)$ per-path cost (H-matrix triangular solve traverses
the block tree), versus the global rSVD's $O(Nk)$ per-path cost.  The
H-matrix advantage lies in construction: $O(Nk\log N)$ vs $O(N^2 k)$ for the
global rSVD sketch, a factor of $N/\log N$ for $k \ll N$.

For the rough case $H = 0.10$, the off-diagonal block needs rank $3$ rather than
$2$ for $1\%$ accuracy (Figure~\ref{fig:structure}(d)): slightly more than at
$H = 0.50$, but far below the ranks the global approximation uses.  The H-matrix
benefit over the global rSVD grows as $N$ increases beyond the current test range.  At $N = 252$,
$O(N^2 k) = O(252^2 \cdot 64) \approx 4 \times 10^6$ flops for construction;
the H-matrix would reduce this by roughly $\log_2 252 \approx 8$ for equal
per-block rank, which becomes worthwhile at $N \geq 1000$.

\subsection{Beyond the two extensions}
\label{sec:ext-future}

Each extension of §\ref{sec:extensions} leaves one gap.  The variance correction gets
the marginal variances right but the increments wrong; replacing the diagonal noise by
a banded factorization of the residual $C - L_k L_k^\top$ (bandwidth $b$) would also
restore the near-diagonal covariances that carry the roughness, at $O(Nb)$ extra cost
per path.  The control variate removes the price-shock noise but leaves the
volatility-path noise, about 4\% of the original variance; a second control built on
the volatility path, for instance the integrated variance $\int_0^T \sigma_t^2\,dt$ whose
mean~\eqref{eq:intvar} is known in closed form, could remove part of what remains.

\subsection{Beyond one machine: GPUs and remaining per-path costs}
\label{sec:gpu-future}

The batched pricers of §\ref{sec:perf} are already in the form a GPU needs: the
Cholesky and low-rank samplers reduce to one dense product per block (cuBLAS
\texttt{trmm}/\texttt{gemm}), the circulant sampler to one batched FFT (cuFFT), and the
counter-style block streams map onto per-thread generators such as Philox.  Two caveats
from the CPU measurements carry over.  Once the transform is batched and the generator is fast, the price
path and its $2N$ exponentials are the largest cost of a low-rank path and about a third
of an FFT path, so the payoff loop would have to be vectorized as well.  And the fair
baseline is the batched, multithreaded CPU code, not the per-path loop, which is
already $1.6$--$3.7\times$ slower on one thread.

\subsection{Hybrid Scheme for rough Bergomi}
\label{sec:hybrid-future}

The Hybrid Scheme of \citet{bennedsen2017hybrid} achieves more than $80\%$
RMSE reduction over naïve Euler discretization at $O(N \log N)$ cost by
exploiting the separability of the Volterra kernel $K(t,s) = (t-s)^{H-1/2}$
near $t = s$: the singular part is handled analytically and the smooth
remainder is convolved via FFT.  This is a structurally superior simulation
method for the rough Bergomi model of \citet{bayer2016pricing}, in
which log-volatility is a Volterra integral $\int_0^t (t-s)^{H-1/2}\,dW_s$.

The RFSV model studied here uses a different representation: $\log\sigma_t =
\nu W_t^H$, where $W_t^H$ is a standard fBM (not a Volterra integral with
the Bergomi kernel).  The two models are different processes with different
correlation structures, though both have Hurst exponent $H$ controlling
roughness.  Adopting the Hybrid Scheme would require switching to the rough
Bergomi model class, which has distinct option pricing properties and
calibration procedures.  The structural insight of separating the kernel into
a singular analytic part and a smooth FFT-able remainder is directly
relevant to the fBM kernel $C(t,s)$ and motivates the hierarchical treatment
in §\ref{sec:hmatrix-future}.

\subsection{Spot-volatility correlation}
\label{sec:rho-future}

The engine draws the price shocks $Z$ independently of the fBM driving
volatility, so the spot-vol correlation is $\rho = 0$.  As §\ref{sec:iv} shows,
this leaves the model unable to produce the downward skew of equity smiles.
Introducing $\rho \neq 0$ requires simulating $W^H$ jointly with the Brownian motion
that drives the price.  The dense Cholesky sampler can absorb this directly by
factoring the $2N \times 2N$ joint covariance, which is explicit when $W^H$ is
written as a Volterra integral of that Brownian motion; the circulant embedding
would need to be extended to the joint process.  This is the single change most
needed before the IV comparison can test the model rather than its
simplifications.

\subsection{Quasi-Monte Carlo path generation}
\label{sec:qmc-future}

All three algorithms use pseudo-random Gaussian draws, giving a standard MC
error of $O(M^{-1/2})$.  Replacing these with low-discrepancy sequences
\citep{niederreiter1992random} reduces the effective error rate to
approximately $O(M^{-1}(\log M)^N)$, close to $O(M^{-1})$ in practice
\citep{glasserman2004monte}.  This improvement is orthogonal to the
choice of covariance simulation method: any of the three algorithms can be
combined with QMC by replacing the standard normal draws.

In $N = 252$ dimensions the low-discrepancy property is strongest for the
first few coordinates.  Brownian bridge dimension ordering concentrates the
largest variance components in the lowest-index dimensions, ensuring the
low-discrepancy sequence acts on the high-information components first
\citep{joy1996quasi}.  At $M = 10{,}000$, the MC standard error of
$\approx 0.095$ (§\ref{sec:convergence}) would fall to $\approx 0.001$ under
ideal $O(M^{-1})$ QMC, a $100\times$ reduction at no additional computational
cost, improving the statistical resolution of the roughness premium and
parameter sensitivity results.
